\documentclass{article}
\usepackage{graphicx}
\usepackage{amsmath}
\usepackage{float}
\usepackage{siunitx}
\usepackage{caption}
\usepackage{subcaption}
\usepackage{hyperref}

\title{Building a Function Generator}
\author{Lily Nguyen}
\date{December 4, 2025}

\begin{document}

\maketitle

% ------------------------
% Introduction and Theory
% ------------------------

\section{Introduction and Theory}

% Each lab report, even if it consists if multiple parts, is centered around a central topic.
% In the introduction, you should introduce this topic and place it in context. When giving
% context, think of answering questions such as, “Why is this technique or tool important?”
% and “Which instruments might rely on this technique or tool?”. Finally, this section should
% include the concepts and equations relevant to your stated goals and the results you plan on
% presenting.

intro



% ------------------------
%   Methods and Results
% ------------------------

\section{Methods and Results}

% Depending on the scientific article (and the style guide of the journal which publishes the
% research), these two sections may or may not be separate. Sometimes, the methods are
% described first and the results come second. Other times the results are shown first and
% the methods are described in a subsequent section. And finally, it sometimes makes sense
% to describe a method and then present the relevant results and then move onto the next
% method. You can choose which of these styles suits each lab report best as long as you
% include all the necessary information in a logical order.

methods/results



% ------------------------
%        Methods
% ------------------------
\subsection{Methods}

% A description of what you did should be presented with enough detail to allow for your
% experiment to be repeated. Imagine that one of your classmates who hasn’t read the lab
% manual walks into lab and is handed your report. Would they be able to reproduce your
% results based on the details you provide in the methods section? You can use words, diagrams,
% pictures (please use figure captions).

methods




% ------------------------
%        Reesults
% ------------------------

\subsection{Results}

% Please make sure to include all the measurements you are asked to perform in the lab manual.
% When presenting the results in a graph, please format the graph with legible axes that are
% titled and labeled with units as well as a figure caption that give all relevant information.

results



% ------------------------
%       Conclusion
% ------------------------

\section{Conclusion}

% Quickly summarize your results and their relevance to electronic methods or techniques. The
% reader should walk away with a clear “take-home message” after reading your conclusion. If
% something went wrong with your experiment and you didn’t get conclusive results, or your
% results disagree with theory, this is your chance to explain what went wrong and propose an
% improvement to the experiment.

conclusion




\end{document}


% ------------------------
%      Style Points
% ------------------------
% Your lab report should flow like a cohesive story. Each next section should be set up by the
% previous one and the reader should never have to ask, “why did they do this thing?”. Run on
% sentences, confusing formatting, badly formatted diagrams and graphs, missing references,
% and excessively long lab reports will result in points deducted in this category.


% ------------------------
%      General Notes
% ------------------------
% There is not a fixed length for any lab report. A lab report is too short if it doesn’t
% flow smoothly because key details required for understanding the material are missing.
% However, please avoid writing a needlessly long lab (the latter is an easier trap to fall
% into). So, please try to be concise!

% Please cite relevant facts and equations unless you have derived them yourself (and unless
% they are considered common knowledge like Ohm’s law). Any citation style is okay as
% long as the source material can be found.
